\chapter{Summary and Outlook}
\label{chap:outlook}

We developed a method for stochastic optimization problems subject to almost sure affine constraints
of the form
\begin{equation*}
    \begin{aligned}
    &\min_{x \in \mathbb{R}^n} \, \{\, f(x) := \E(F_\xi(x)) + r(x) \,\} \\
    &\textup{subject to} \quad A(\xi) x - b(\xi) \leq 0 \quad \textnormal{almost surely}.
    \end{aligned}
\end{equation*}
Our method naturally handles the case where there are infinitely many constraints, which is more general than
than the settings which have been considered in related works \cite{nedich2023huber,li2025new}. For example,
our guarantees hold in the setting of
robust optimization by sampling
\cite{ben2009robust,xu2020primal,calafiore2005uncertain}.
Further, because our algorithm
uses penalty functions to softly enforce constraints, it is applicable even if the constraints are not fully satisfiable.
Empirically, we found that the iterates still converge to a reasonable solution in that case (see \cref{sec:numerical-example-svm}).

We analyzed the convergence
properties of the method in the convex and strongly convex regimes.
In the convex setting, one can achieve a convergence rate of
$\mathcal{O}(\log^{1/2 + \epsilon} K / \sqrt{K})$ for any $\epsilon \in (0,\infty)$.
The convergence rate can be improved
to $\mathcal{O}(\log^{\epsilon} K / K)$ for any $\epsilon \in (0,\infty)$, in the strongly convex case.
These rates improve upon the rates achieved in the closely related work by Fercoq et al. \cite{fercoq2019almost}
without the need of a global variance bound on the stochastic gradients.
Furthermore, it is known that stochastic projected gradient descent achieves the minimax optimal convergence rates of
$\mathcal{O}(1/\sqrt{K})$ in the convex setting, and $\mathcal{O}(1/K)$ in the strongly convex setting \cite{nemirovski_yudin_1983,agarwal2009information}.
Thus, the convergence rates of our method match the optimal convergence rates
up to an arbitrarily small logarithmic factor, while being projection-free.
Besides the typical convergence guarantees which hold in expectation, we also provided convergence guarantees that hold
with high-probability in the general convex setting.
Notably, we did not rely on any of the typical
global boundedness assumptions regarding the stochastic gradients of the objective $f$.

There are a lot of research directions for future work.
The most glaring restriction is that we only handle constraints that
are affine in the decision variable.
An extension to more general functional inequalities would be a significant next step.
It would also be interesting to analyze the algorithm in the more common setting of chance constraints, where
constraints only need to hold with some
high probability, instead of almost surely.
A straight-forward way to transform a chance constrained problem into an
almost sure one is through the sampling-and-discarding approach
\cite{campi2011sampling}. However, this requires the ability to sample
a possibly very large amount of times, which is not feasible
in \textit{online optimization} \cite{shalev2025online} settings, where one typically cannot sample on-demand.

Another possible avenue for future research is the use of acceleration methods.
To speed up the convergence of first-order methods, a popular method in deterministic optimization
is Nesterov's \textit{accelerated
gradient method} (AGD) \cite{nesterov1983method,nesterov2004introductory}. In stochastic optimization, a naive application
of AGD does not lead to faster (asymptotic) convergence rates, due to an increase in variance in the gradient estimates.
Instead, acceleration is achieved via variance reduction methods, such as SAGA \cite{defazio2014saga} and SVRG \cite{johnson2013svrg}.
These algorithms achieve acceleration, while keeping variance managable by periodically calculating full gradients of the objective.
Unfortunately, a straight-forward application of these methods to our setting is not possible,
unless we assume a finite-sum structure of the objective and a finite number of constraints.
Moreover, the number of summands involved in the objective, as well as the number of constraints, need to be managable, so
that the occasional calculation of full gradients of the penalized objective used in our method is feasible.
A promising recent method that circumvents these issues was proposed in \cite{lei2017less}, where the authors manage
to greatly improve the convergence of stochastic gradient methods, while avoiding the calculation of full gradients.
An adaptation of their methods to our constrained setting is a promising future research direction to improve upon the convergence
rate of our algorithm.
